\emph{(1) The radical pairing.} Let $(\widetilde W, Q)$ be any completion. For $\rho \in R$ the functional $Q(\rho, \cdot)$ vanishes on all of $W$ (it equals $B(\rho, \cdot)$ there), hence descends to a functional on $\widetilde W / W$. This defines a linear map
\[
\lambda : R \longrightarrow (\widetilde W / W)^*, \qquad \lambda(\rho) := Q(\rho, \cdot) \bmod W .
\]
If $\lambda(\rho) = 0$ then $Q(\rho, \cdot)$ vanishes on $W$ and on representatives of $\widetilde W / W$, hence on $\widetilde W$, so $\rho = 0$ by nondegeneracy: $\lambda$ is injective, and $\dim(\widetilde W / W) \geq r$.

\emph{(2) Existence at the bound, and the forced signature.} Take $\widetilde W := W \oplus R^*$ and define $Q$ by: $Q|_W = B$; $Q(\rho, \varphi) = \varphi(\rho)$ for $\rho \in R$, $\varphi \in R^*$; $Q(\varphi, \varphi') = 0$ on $R^*$; and $Q(V_1, R^*) = 0$ for a fixed complement $V_1$ of $R$ in $W$. Nondegeneracy: if $v_1 + \rho + \varphi$ ($v_1 \in V_1$, $\rho \in R$, $\varphi \in R^*$) is $Q$-orthogonal to everything, then pairing against $R^*$ gives $\rho = 0$, against $R$ gives $\varphi = 0$, and against $V_1$ gives $v_1 = 0$ since $B|_{V_1}$ is nondegenerate (it is a complement of the radical). For the signature, choose a basis $\rho_1, \dots, \rho_r$ of $R$ with dual basis $\varphi_1, \dots, \varphi_r$: each $\operatorname{span}\{\rho_i, \varphi_i\}$ is a hyperbolic plane, the planes are mutually orthogonal and orthogonal to $V_1$, so $Q \cong B|_{V_1} \perp H^{\perp r}$ has signature $(p{+}r, q{+}r)$.

Now let $(\widetilde W, Q)$ be \emph{any} completion of dimension $m + r$. Then $V_1^{\perp}$ (orthocomplement in $\widetilde W$) is nondegenerate of dimension $(m{+}r) - (p{+}q) = 2r$, being the orthocomplement of a nondegenerate subspace, and it contains $R$ (radical vectors are $B$-orthogonal to $V_1$), which is totally isotropic of dimension $r$: half the dimension. We claim:
\begin{quote}
\emph{a nondegenerate quadratic space $(X, Q_X)$ of dimension $2s$ containing a totally isotropic subspace $U$ of dimension $s$ is isometric to $H^{\perp s}$.}
\end{quote}
The claim with $s = r$ gives $Q|_{V_1^\perp} \cong H^{\perp r}$ of signature $(r,r)$, and the total is $(p{+}r, q{+}r)$.

Proof of the claim, by induction on $s$, the case $s = 0$ being empty. Pick $u \in U$ nonzero; by nondegeneracy there is $y$ with $Q_X(u, y) = 1$, and $y' := y - \tfrac12 Q_X(y,y)\, u$ is null with $Q_X(u, y') = 1$ still, so $H_1 := \operatorname{span}\{u, y'\}$ is a hyperbolic plane and $X = H_1 \perp X_1$ with $X_1 := H_1^{\perp}$ nondegenerate of dimension $2s - 2$. For $u' \in U$ set $u'' := u' - Q_X(u', y')\, u$: then $u''$ pairs to zero with $u$ (both terms do, $U$ being isotropic) and with $y'$ (the correction cancels), so $u'' \in X_1$. The corrected vectors span a totally isotropic subspace of $X_1$, and the correction map $u' \mapsto u''$ is linear with kernel exactly $\R u$, so that subspace has dimension $s - 1$. Induction finishes the claim.

\emph{(3) Uniqueness.} Let $(\widetilde W_1, Q_1)$ and $(\widetilde W_2, Q_2)$ be minimal completions. By (2) they are isometric abstractly; more precisely, the identity $W \to W$ is an isometry between the subspaces $W \subseteq \widetilde W_1$ and $W \subseteq \widetilde W_2$ of two isometric nondegenerate spaces, and by Fact \claimref{conformal.witt-extension-theorem} it extends to an isometry $\widetilde W_1 \to \widetilde W_2$. Functoriality of $\Cl$ turns this isometry into an algebra isomorphism fixing the embedded copy of $\Cl(p,q,r)$.
